\documentclass[10pt,a4paper]{amsart}
\usepackage[margin=19mm]{geometry}
\usepackage{amsmath,amssymb,mathtools}
\usepackage[scr=rsfs,cal=euler]{mathalfa}
\usepackage{xfrac}
\usepackage[unicode,hidelinks,bookmarks=false]{hyperref}
\newcommand{\ie}{i.e.\ }
\newcommand{\cf}{cf.\ }
\newcommand{\resp}{respectively\ }
\newcommand{\Subsection}{\subsection}
\providecommand{\nobreakdash}{\nobreak\hbox{-}}
\setlength{\emergencystretch}{2em}
\multlinegap0pt
\usepackage{bm}
\usepackage{bbm}
\usepackage{dsfont}
\usepackage{xspace}
\usepackage{cancel}
\usepackage{mathtools}
\usepackage{amsthm}
\makeatletter
\@ifundefined{thm}{\newtheorem{thm}{Thm}}{}
\@ifundefined{cor}{\newtheorem{cor}[thm]{Corollary}}{}
\makeatother
\def\a{\alpha}
\def\d{\delta}
\title[Topics in Magnetic Geometry: Interpolation, Intersections an]{Topics in Magnetic Geometry: Interpolation, Intersections and Integrability}
\author{Lina Deschamps, Levin Maier, Tom Stalljohann}
\date{Source: arXiv:2604.13616v2 (CC BY 4.0)}
\begin{document}
\maketitle

\noindent\emph{Source and licence.} Topics in Magnetic Geometry: Interpolation, Intersections and Integrability. Version arXiv:2604.13616v2 (CC BY 4.0), distributed under the Creative Commons Attribution 4.0 International licence, as recorded in the arXiv licence metadata for this exact version and shown on the abstract page https://arxiv.org/abs/2604.13616v2. Full rights notice: licenses/arxiv-2604.13616-CC-BY-4.0.txt.\par\smallskip

\noindent\textit{Editorial scope.} Counted unit: the Cor whose source label is \texttt{cor: total integrability in dim 3} in the article (source0.tex lines 986--1010, source label \texttt{cor: total integrability in dim 3}), delivered together with the article's own complete original proof of it (source0.tex lines 1012--1089, one \texttt{proof} environment of 78 lines). No mathematics is written, repaired or re-derived: statement and proof are the article's own text line for line. Required context retained uncounted: it:prop hyp (lines 827--830); eq: derivative of contact angle (lines 852--856). Every definition, display and label of those lines is retained unchanged. Omitted from this package and not redistributed: 1--826, 831--851, 857--985, 1011--1011, 1090--1832. Nothing inside a retained excerpt is omitted or altered: every retained body is delivered exactly as the article's lines are written, so no omission receipt of any kind accompanies this package. Citations: the retained excerpts carry no citation command at all, so no bibliography entry of the article is transcribed and no reference is redistributed. Reference adaptation: none was needed and none was made; every reference inside the delivered unit resolves inside it, so no unresolved reference or citation is produced. Printed slips kept literal: no printed word, symbol, hypothesis or number of the counted statement or of its proof was altered. Layout and class: the article's own class is \texttt{amsart} with packages including inputenc, a4wide, appendix, graphicx, enumerate, hyperref, cleveref, amsmath, amsopn, amssymb, amsfonts, stmaryrd, verbatim, todonotes; the wrapper is the lane's pinned \texttt{amsart} editorial wrapper at 10pt on A4, which restates the theorem-like environments the retained text needs and adds the article's own macro definitions that the retained text needs (\texttt{\textbackslash a}, \texttt{\textbackslash d}), with extra wrapper packages (bm, bbm, dsfont, xspace, cancel, mathtools). The delivered numbering is the unit's own. Assets: no retained span contains a figure environment, a graphics-inclusion command, a PDF-page inclusion, a table or a TikZ picture; no image file is copied into this package, the package declares no asset dependency and no image file is redistributed with it. Licence: version arXiv:2604.13616v2 of this article is distributed under the Creative Commons Attribution 4.0 International licence, as recorded in the arXiv licence metadata for this exact version and shown on the abstract page https://arxiv.org/abs/2604.13616v2. A full rights notice is delivered at licenses/arxiv-2604.13616-CC-BY-4.0.txt. Characters: every retained excerpt is pure ASCII, so the delivered engine, XeLaTeX with its default Latin Modern text font, reports no missing character.
\begin{equation}\label{it:prop hyp}
    \d\alpha(X,\cdot)=0 \, ,
%    X_p\in \ker \d\alpha_p \quad \quad \text{for all } \ p\in M\,
\end{equation}

\begin{equation}\label{eq: derivative of contact angle}
    \frac{\mathrm{d}}{\mathrm{d} t} g_{\gamma}(X_{\gamma}, \dot{\gamma}) 
    = g_{\gamma}\left(\nabla_{\dot{\gamma}} X_{\gamma}, \dot{\gamma} \right) 
    + g_{\gamma}\left(X_{\gamma}, Y_{\gamma}(\dot{\gamma}) \right).
\end{equation}

\begin{cor}\label{cor: total integrability in dim 3}
Let \( (M^3,g,\d\alpha) \) be a Killing magnetic system on a closed
three-dimensional manifold \( M^3 \) such that $\d\a\neq 0$ and \eqref{it:prop hyp} holds. Then the functions 

\[
H_1(p,v):=\tfrac12 |v|^2,\qquad
H_2(p,v):=g_p(X_p,v),\qquad
H_3(p,v):=g_p(X_p,X_p)=|X|^2(p).
\]
are poisson commuting integrals of motion of the magnetic geodesic flow $\varPhi_{g,\d\a}^t$. 
Moreover, these three first integrals are functionally independent at every
\((p,v)\in TM\) such that
\[
v\notin \mathbb R X_p
\qquad\text{and}\qquad
\d(|X|^2)_p\neq 0.
\]
Consequently, the magnetic geodesic flow is Liouville integrable on every
connected component of the open set
\[
\mathcal R
:=
\{(p,v)\in TM:\ v\notin \mathbb R X_p,\ \d(|X|^2)_p\neq 0\}.
\]
\end{cor}

\begin{proof}
By \eqref{eq: derivative of contact angle}, the function \(H_2\) is a first integral of the magnetic geodesic flow \( \varPhi^t_{g,\mathrm d\alpha} \). By \eqref{it:prop hyp}, the function \(H_3\) is also a first integral. Therefore \[ \{H_1,H_2\}=0,\qquad \{H_1,H_3\}=0. \]
Now \(H_2\) is the standard momentum map associated with the Killing field
\(X\), hence its Hamiltonian flow is the tangent lift of the flow of \(X\).
Since \(X\) is Killing, its flow preserves \(X\) itself, and therefore also
preserves its squared norm. Thus
\[
X(H_3)=X(|X|^2)=0.
\]
Equivalently, \(H_3\) is invariant under the Hamiltonian flow generated by
\(H_2\), and so
\[
\{H_2,H_3\}=0.
\]
Therefore \(H_1,H_2,H_3\) are pairwise Poisson--commuting first integrals of
the magnetic geodesic flow.

It remains to prove the functional independence statement. Fix
\((p,v)\in TM\). Using the canonical splitting
\[
T_{(p,v)}TM \simeq T_pM \oplus T_pM,
\]
write a tangent vector as \((\xi,\eta)\), where \(\xi\) is the horizontal part
and \(\eta\) the vertical part. Then
\[
\d H_1(\xi,\eta)=g_p(v,\eta),
\]
\[
\d H_2(\xi,\eta)=g_p(\nabla_\xi X,v)+g_p(X_p,\eta),
\]
and
\[
\d H_3(\xi,\eta)=\xi(|X|^2).
\]
Suppose that at \((p,v)\) one has a linear relation
\[
a\,\d H_1+b\,\d H_2+c\,\d H_3=0 .
\]
Evaluating this on vertical vectors \((0,\eta)\), we obtain
\[
a\,g_p(v,\eta)+b\,g_p(X_p,\eta)=0
\qquad\text{for all }\eta\in T_pM.
\]
Hence
\[
a\,v+b\,X_p=0.
\]
If \(v\notin \mathbb RX_p\), this implies
\[
a=b=0.
\]
Therefore the above relation reduces to
\[
c\,\d H_3=0.
\]
If in addition \(\d(|X|^2)_p\neq 0\), then \(\d H_3\neq 0\) at \((p,v)\), so
\(c=0\). Thus
\[
\d H_1,\ \d H_2,\ \d H_3
\]
are linearly independent at every \((p,v)\in TM\) such that
\[
v\notin \mathbb RX_p
\qquad\text{and}\qquad
\d(|X|^2)_p\neq 0.
\]

Consequently, on every connected component of the open set
\[
\mathcal R
=
\{(p,v)\in TM:\ v\notin \mathbb RX_p,\ \d(|X|^2)_p\neq 0\},
\]
the magnetic geodesic flow on the \(6\)-dimensional symplectic manifold \(TM\)
admits three functionally independent Poisson--commuting first integrals
\(H_1,H_2,H_3\). Hence, by the Liouville--Arnold theorem, the flow is
Liouville integrable on every connected component of \(\mathcal R\).
\end{proof}

\end{document}
